
\documentclass[12pt,a4paper]{report}

% Configuração do idioma e codificação
\usepackage[utf8]{inputenc}
\usepackage[T1]{fontenc}
\usepackage[portuguese]{babel}

% Margens e formatação
\usepackage[a4paper,margin=2.5cm]{geometry}
\usepackage{graphicx}
\usepackage{amsmath}
\usepackage{hyperref}
\usepackage{setspace}

\onehalfspacing

% Informações do documento
\title{Sistema Inteligente de Apoio à Decisão
de Treinadores de Basquetebol}

\author{Relatório de AAUTIA e PLANTIA}

\date{\today}

\begin{document}

\maketitle

\begin{abstract}
Este trabalho investiga a integração de técnicas
de Machine Learning e Planeamento Automático
no apoio à decisão de treinadores de basquetebol.
\end{abstract}

\tableofcontents
\newpage

\section{Introdução}

A Inteligência Artificial tem vindo a assumir
um papel crescente na análise de dados desportivos.

O presente trabalho pretende desenvolver um
sistema inteligente de apoio à decisão.

\section{Estado da Arte}

\subsection{Aprendizagem Automática}

A aprendizagem automática permite identificar
padrões nos dados históricos e prever resultados
de jogos de basquetebol.

\subsection{Planeamento em Inteligência Artificial}

O planeamento automático permite selecionar
ações que contribuam para atingir objetivos
competitivos, respeitando restrições.

\section{Metodologia}

Serão utilizados algoritmos de Machine Learning
e métodos de planeamento automático.

\section{Conclusões}

O objetivo é integrar previsão e planeamento
num sistema de apoio às decisões do treinador.

\end{document}
